\documentclass[11pt]{article}
\usepackage[T1]{fontenc}
\usepackage{lmodern}
\usepackage[margin=1in]{geometry}
\usepackage{amsmath,amssymb,amsthm,mathtools}
\usepackage{microtype}
\usepackage[hidelinks]{hyperref}
\usepackage{url}
\hypersetup{pdftitle={Deterministic Connectivity: Audited Baseline and Collision Barriers},
 pdfsubject={Independently prepared research addendum}}
\newtheorem{theorem}{Theorem}[section]
\newtheorem{lemma}[theorem]{Lemma}
\newtheorem{proposition}[theorem]{Proposition}
\theoremstyle{remark}
\newtheorem{remark}[theorem]{Remark}
\newcommand{\F}{\mathbb F_2}
\newcommand{\rank}{\operatorname{rank}}
\newcommand{\Cut}{\mathcal C}
\newcommand{\ceil}[1]{\left\lceil #1\right\rceil}
\title{Deterministic Connectivity:\\Audited Baseline and Collision Barriers}
\author{}
\date{Research addendum, 29 October 2026}

\begin{document}
\maketitle
\begin{abstract}
We give a self-contained proof of the supplied draft's deterministic
$\Omega(\sqrt n)$ lower bound for exact Boolean connectivity in the
simultaneous vertex-neighborhood model, with explicit alphabet and bit
constants. We also establish limits on several proposed collision arguments:
local linear maps with $r+\ceil{\log_2 n}$ bits exclude every nonzero
adjacency difference of binary rank at most $r$; local XOR-linear cut-space
recognizers have concatenated rank at least $\binom m2-1$; a four-vertex
example defeats a particular anchored collision template; and public
$r$-sheet lifts have a direct reconstruction protocol. These results do not
improve the unrestricted $\Omega(\sqrt n)$ lower bound. This is an
independent addendum, with no claim of novelty or author acceptance.
\end{abstract}

\section{Definitions and scope}

An input is a simple undirected graph $G$ on the public labeled vertex set
$[n]$. Player $v$ knows exactly its own neighborhood $N_G(v)$ and sends one
deterministic message to a referee. All messages are simultaneous. Encoders
may be vertex-specific and nonlinear, and the referee has no computational
restriction. The required output is the single bit indicating whether $G$
is connected, with correctness on every input. No shared knowledge of other
vertices' neighborhoods is available.

The communication cost $B$ is the worst-case \emph{maximum bits sent by one
vertex}, not the total number of bits. We distinguish a message alphabet of
cardinality at most $2^b$, fixed-length binary messages, and arbitrary binary
strings of maximum length $B$. All logarithms are base two unless written
$\ln$. Only ordinary connectivity, or $k=1$, is treated here.

For algebraic statements, edge indicators lie in
$\F^{\binom n2}$, with the fixed ambient edge set $\binom{[n]}2$.
An adjacency difference $D$ is a symmetric zero-diagonal matrix over $\F$;
$D_v$ denotes its incident row with the diagonal coordinate omitted.
The ranks of these adjacency matrices and the ranks of concatenated
message maps are different quantities.

\section{The unrestricted connectivity baseline}

\begin{theorem}[D-01, ordinary connectivity]\label{thm:baseline}
Let $n\ge2$ and let $b\ge1$ be an integer. If an exact deterministic
connectivity protocol has at most $2^b$ possible messages at each vertex,
then
\begin{equation}\label{eq:baseline}
 b>\frac{n}{768(\ceil{\sqrt n}+1)}
   >\frac{\sqrt n}{2304}.
\end{equation}
Consequently, fixed-length messages of maximum length $B$ satisfy
$B>\sqrt n/2304$. Variable-length binary strings of maximum length $B$,
including the empty string, satisfy $B>\sqrt n/4608$.
\end{theorem}

We first prove the counting fact used to select a common partition.
A family $\mathcal S\subseteq\{1,2,3,4\}^q$ \emph{shatters} a coordinate
set if its projection onto those coordinates contains every assignment.

\begin{lemma}[Four-symbol shattering bound]\label{lem:shattering}
If $\mathcal S\subseteq\{1,2,3,4\}^q$ shatters no coordinate set of
size $d$, then
\[
 |\mathcal S|\le F(q,d):=
 \sum_{j=0}^{d-1}\binom qj3^{q-j}.
\]
Here $F(q,0)=0$ and $F(q,d)=4^q$ for $d>q$.
\end{lemma}

\begin{proof}
The cases $d=0$ and $d>q$ are immediate, since any nonempty family shatters
the empty coordinate set. Induct on $q$, with $q=0$ immediate. Let
$\mathcal S_a$ be the slice having last coordinate $a$, projected onto
the first $q-1$ coordinates. Set
$\mathcal U=\bigcup_{a=1}^4\mathcal S_a$ and
$\mathcal I=\bigcap_{a=1}^4\mathcal S_a$. Each projected word occurs
in at most three slices unless it is in $\mathcal I$, so
\[
 |\mathcal S|=\sum_{a=1}^4|\mathcal S_a|
 \le3|\mathcal U|+|\mathcal I|.
\]
The family $\mathcal U$ cannot shatter $d$ coordinates. The family
$\mathcal I$ cannot shatter $d-1$ coordinates, since adjoining the last
coordinate would then give a shattered $d$-set in $\mathcal S$.
Induction and Pascal's identity give
\[
 |\mathcal S|\le3F(q-1,d)+F(q-1,d-1)=F(q,d).
\]
\end{proof}

\begin{proof}[Proof of Theorem~\ref{thm:baseline}]
Fix a correct protocol and put
\[
 q=128(\ceil{\sqrt n}+1),\qquad M=n-q.
\]
If $b>\sqrt n$, the asserted alphabet bound is immediate, so assume
$b\le\sqrt n$. If $M<4(qb+1)$, then, using $q\ge4$ and $b\ge1$,
\begin{equation}\label{eq:bookkeeping}
 n<(4b+1)q+4\le(4b+2)q\le6bq,
\end{equation}
as desired. We therefore suppose, towards a contradiction, that
\begin{equation}\label{eq:largeM}
 M\ge4(qb+1).
\end{equation}
This implies $q<n$. Choose a reservoir $U$ of $q$ vertices and designate
the remaining $M$ vertices as centers $c_1,\ldots,c_M$. For each center
let $f_i(S)$ be its actual message when its neighborhood is $S\subseteq U$.

\smallskip
\noindent\emph{Local witnesses.}
Set $d=16(b+1)$. Consider any $d$-element set $Y\subseteq U$ and one
center $c_i$. Pair the vertices of $Y$ and designate a pair
$\{\alpha,\beta\}$. Consider all sets $A$ containing $\alpha$, excluding
$\beta$, and choosing exactly one vertex of each other pair. There are
$2^{d/2-1}=2^{8b+7}>2^{2b}$ such sets. The ordered message pair
$(f_i(A),f_i(Y\setminus A))$ has at most $2^{2b}$ values. Thus distinct
sets $A_0,A_1$ have the same ordered pair of messages. Define
\[
\begin{array}{ll}
 P_i^0=A_0\cap A_1,&P_i^1=Y\setminus(A_0\cup A_1),\\
 R_i^0=A_0\setminus A_1,&R_i^1=A_1\setminus A_0.
\end{array}
\]
These sets partition $Y$. The first two contain $\alpha$ and $\beta$,
respectively. The last two are nonempty because the two choices differ on
at least one entire pair. Equality of the two ordered messages yields
\begin{equation}\label{eq:witness}
\begin{split}
 f_i(P_i^0\cup R_i^0)&=f_i(P_i^0\cup R_i^1),\\
 f_i(P_i^1\cup R_i^0)&=f_i(P_i^1\cup R_i^1).
\end{split}
\end{equation}
Label these four sets by $P_0,P_1,R_0,R_1$, respectively. Every extension
of this labeling to $U$ contains the same nonempty witnessing subsets.

\smallskip
\noindent\emph{One partition for every center.}
Call a labeling $U\to\{P_0,P_1,R_0,R_1\}$ good for $c_i$ if it admits
four nonempty subsets in the indicated classes satisfying
\eqref{eq:witness}, and bad otherwise. The local argument shows that the
bad family shatters no $d$-element coordinate set. Under an independently
uniform labeling of $U$, Lemma~\ref{lem:shattering} gives
\[
 \Pr[\text{bad for }c_i]
 \le4^{-q}\sum_{j<d}\binom qj3^{q-j}
 =\Pr[Z<d],\qquad Z\sim\operatorname{Bin}(q,1/4).
\]
Since $b\le\sqrt n$, we have $d\le q/8$. Write $\mu=q/4$ and
$t=\ln2$. Exponential Markov's inequality implies
\begin{align*}
 \Pr[Z\le\mu/2]
 &\le e^{t\mu/2}\mathbb E e^{-tZ}
 =e^{t\mu/2}(3/4+e^{-t}/4)^q\\
 &\le\exp\{\mu(t/2-1/2)\}
 \le e^{-\mu/8}=e^{-q/32}.
\end{align*}
Here $1+x\le e^x$ and $(\ln2)/2-1/2<-1/8$ justify the last two
steps. A union bound over the $M\le n$ centers gives
\[
 \Pr[\text{some center is bad}]
 \le ne^{-q/32}\le ne^{-4\sqrt n}<1.
\]
For the final inequality, setting $x=\sqrt n\ge1$ gives
$2\ln x-4x<0$. Therefore a partition
$U=P_0\sqcup P_1\sqcup R_0\sqcup R_1$ is good for every center.
Fix it and fix witnessing subsets for every center. No independence of
the centers' bad events was used.

\smallskip
\noindent\emph{A global transcript collision.}
For $a,z\in\{0,1\}^M$, define $G(a,z)$ by making $P_0\cup R_0$ and
$P_1\cup R_1$ into cliques and giving $c_i$ exactly the neighborhood
$P_i^{a_i}\cup R_i^{z_i}$. There are no other edges. Both cliques are
nonempty. The graph $G(x,x)$ has exactly two components, because every
center attaches within just one clique. If $x\ne y$, a center at a
differing coordinate joins the two cliques in $G(x,y)$; all other
centers have a neighbor in a clique. Thus $G(x,y)$ is connected.

By \eqref{eq:witness}, each center's message depends only on $a_i$.
A vertex in $P_0\cup P_1$ sees its fixed clique and centers selected
only by $a$, so its message depends only on $a$. A vertex in
$R_0\cup R_1$ analogously has a message depending only on $z$.
Their joint transcript $T(z)$ has at most $2^{qb}$ values. Since
\eqref{eq:largeM} implies $M>qb$, there are distinct $x,y$ with
$T(x)=T(y)$. The full transcripts of $G(x,x)$ and $G(x,y)$ agree:
the center and $P$ messages agree because $a=x$ is fixed, and the $R$
messages agree by construction. Their connectivity differs, a
contradiction.

It follows that $M<4(qb+1)$, so \eqref{eq:bookkeeping} proves
$b>n/(6q)$. This argument also covers all small $n$: when $q\ge n$,
the inequality $M<4(qb+1)$ already holds and no reservoir is constructed.
Finally $\ceil{\sqrt n}+1<3\sqrt n$ for $n\ge2$, which proves the
second inequality in \eqref{eq:baseline}.

Zero-bit protocols are impossible, since the empty and complete graphs
require different outputs. Fixed-length messages therefore have $B\ge1$,
and we apply the alphabet result with $b=B$. For variable-length strings,
the available alphabet has size
$\sum_{j=0}^B2^j=2^{B+1}-1<2^{B+1}$. Apply the result with $b=B+1$ and
use $B+1\le2B$ to obtain $B>\sqrt n/4608$.
\end{proof}

\section{A barrier to bounded-rank collision witnesses}

\begin{theorem}\label{thm:barrier}
For integers $n\ge2$ and $r\ge0$, there are vertex-specific linear maps
\[
 A_v:\F^{n-1}\longrightarrow\F^b,
 \qquad b=r+\ceil{\log_2 n},
\]
such that every nonzero symmetric zero-diagonal $n\times n$ matrix $D$
of rank at most $r$ satisfies $A_v(D_v)\ne0$ for some vertex $v$.
Consequently, two graph inputs with a nonzero adjacency-XOR difference
of rank at most $r$ have distinct concatenated transcripts under these maps.
\end{theorem}

The maps in this theorem are deterministic after being fixed by an
existence argument. The theorem supplies neither a connectivity-deciding
referee nor an efficient uniform construction of the maps. A chosen dense
map can be evaluated in $O(nb)$ binary operations per vertex.

\begin{lemma}\label{lem:alternating}
For $t\ge1$ and $r\ge0$, there are at most $2^{rt}$ symmetric
zero-diagonal $t\times t$ matrices over $\F$ of rank at most $r$.
\end{lemma}

\begin{proof}
Such a matrix represents an alternating bilinear form. If the form is
nonzero, choose two vectors with pairing one. Their two-dimensional span
is nondegenerate and splits off from its orthogonal complement. Iterating
gives a basis in which the form is a direct sum of blocks
$\left(\begin{smallmatrix}0&1\\1&0\end{smallmatrix}\right)$ and a zero
block. Thus the rank is even, say $2k$, and the matrix can be written
$U J_{2k}U^{\mathsf T}$, where $J_{2k}$ is the sum of $k$ such blocks
and $U$ has $t$ rows and $2k$ columns. This follows by transforming the
displayed normal form back to the original basis.

Put $r'=2\lfloor r/2\rfloor$. Padding $U$ with zero columns represents
every matrix of rank at most $r$ as $U J_{r'}U^{\mathsf T}$ with
$U\in\F^{t\times r'}$. There are $2^{r't}\le2^{rt}$ choices of $U$.
The representation need not be unique. If $r'=0$, it represents only
the zero matrix, as required.
\end{proof}

\begin{proof}[Proof of Theorem~\ref{thm:barrier}]
Choose every entry of every local matrix $A_v$ independently and
uniformly in $\F$. Fix a nonzero adjacency difference $D$ and let $T$
be its set of nonzero rows, with $t=|T|$. Each row $D_v$, $v\in T$,
is nonzero, so its image under $A_v$ is uniformly distributed in $\F^b$.
The local maps are independent. Hence
\[
 \Pr[A_v(D_v)=0\text{ for all }v]=2^{-bt}.
\]
Symmetry confines every nonzero entry to the principal submatrix on $T$.
Moreover $t\ge2$, because a nonzero off-diagonal entry makes both its
rows nonzero. Lemma~\ref{lem:alternating} bounds the number of possible
matrices on each fixed $T$ by $2^{rt}$. Union-bounding over actual row
supports gives
\begin{align*}
 \Pr[\text{some nonzero }D\text{ of rank }\le r\text{ is killed}]
 &\le\sum_{t=2}^n\binom nt2^{(r-b)t}\\
 &\le\sum_{t=2}^n\binom nt n^{-t}
 \le\sum_{t=2}^n\frac1{t!}<1.
\end{align*}
For the strict inequality, $t!\ge2^{t-1}$ gives the upper bound
$\sum_{t=2}^n2^{-(t-1)}=1-2^{-(n-1)}<1$.
The count includes matrices whose support is smaller than $T$, which
only enlarges the bound; the probability $2^{-bt}$ is assigned only
to matrices with actual support $T$. Thus a successful choice of maps
exists. For $r=0,1$ the claim is vacuous, since a nonzero alternating
matrix has rank at least two.
\end{proof}

\begin{remark}
In particular, rank-two differences can all be excluded by
$\ceil{\log_2 n}+2$ bits per vertex. A universal collision claim using
only the message budget therefore cannot require such differences for
every encoder collection. This does not refute an argument using
additional consequences of an assumed correct connectivity protocol,
exclude high-rank collisions, or give an unrestricted connectivity
upper bound. If $b>n-1$, full-row transmission already provides a simpler
separator.
\end{remark}

\section{Two cut-space facts with restricted conclusions}

On the complete graph $K_m$, define the binary cut space
\[
 \Cut_m=\{\delta s:s\in\F^m\},\qquad
 (\delta s)_{uv}=s_u+s_v.
\]
In the local recognition problem, vertex $v$ knows only the edge-label
vector $x_v$ on its star and sends a message; the referee must decide
whether $x\in\Cut_m$. This is a distinct task from unrestricted graph
connectivity.

\begin{theorem}[Local XOR-linear cut recognition]\label{thm:cutrank}
Let $m\ge3$. Suppose local linear maps
$L_v:\F^{m-1}\to\F^{B_v}$, concatenated as $L$, permit an arbitrary
referee to recognize $\Cut_m$ exactly. Then
\[
 \rank L\ge\binom m2-1.
\]
For the maximum fixed-length message budget $B=\max_v B_v$, this gives
\[
 \sum_v B_v\ge\binom m2-1,
 \qquad B\ge\ceil{\frac{\binom m2-1}{m}}.
\]
\end{theorem}

\begin{proof}
Let $K=\ker L$. Every vector in $K$ has the same transcript as the zero
cut, so correctness implies $K\subseteq\Cut_m$. Suppose $\dim K\ge2$,
and choose independent cuts $\delta s,\delta t\in K$. Define
\[
 x_{uv}=s_vt_u+t_vs_u\qquad(u\ne v).
\]
This is a symmetric edge labeling. At vertex $v$,
\[
 x_v=s_v(\delta t)_v+t_v(\delta s)_v,
\]
since the two copies of $s_vt_v$ cancel in characteristic two.
Local linearity gives $L_v(x_v)=0$ at every vertex, hence $x\in K$.

The labels $z_v=(s_v,t_v)\in\F^2$ affinely span $\F^2$. Otherwise
some nonzero $(a,b)$ and constant $c$ satisfy $as_v+bt_v=c$ for every
$v$, implying $a\delta s+b\delta t=0$, contrary to independence.
Choose $p,q,r$ such that $z_q+z_p$ and $z_r+z_p$ are independent.
Direct expansion then gives
\[
 x_{pq}+x_{qr}+x_{rp}
 =\det(z_q+z_p,z_r+z_p)=1.
\]
Every cut has triangle sum zero, so $x\notin\Cut_m$, contradicting
$x\in K\subseteq\Cut_m$. Thus $\dim K\le1$, and rank-nullity gives
the rank bound. Finally $\rank L\le\sum_v B_v\le mB$ proves the
communication consequences.
\end{proof}

The decisive step uses locality and linearity of each encoder.
The theorem imposes no linearity assumption on the referee, but its
proof does not extend to arbitrary nonlinear encoders. Its primary
conclusion concerns the total rank of $L$, rather than a rank bound
at each vertex.

\begin{proposition}[A four-vertex obstruction to an anchored template]
\label{prop:anchor}
Let $S=\{s\in\F^4:s_0=0\}$ and define
\[
 h_0(s)=s_1,\qquad h_1(s)=s_2,\qquad
 h_2(s)=s_3,\qquad h_3(s)=s_1.
\]
There is no $s\in S$ and nonzero symmetric zero-diagonal matrix $D$
with root row $D_0=0$ for which
\begin{equation}\label{eq:anchor}
 h_v(s+D_v)=h_v(s)\quad\text{for every }v.
\end{equation}
Here $D_v$ is the full four-coordinate row, including its zero diagonal
entry. Nevertheless, the associated local messages
$h_v(P_v(x))$, where
\[
 P_v(x)_i=x_{vi}+x_{v0}\quad(0\le i\le3),\qquad x_{vv}=0,
\]
do not recognize $\Cut_4$.
\end{proposition}

\begin{proof}
The anchor condition makes $D_{01}=D_{02}=D_{03}=0$, and also ensures
that every $s+D_v$ belongs to $S$. Since the $h_v$ are linear,
\eqref{eq:anchor} at vertices $1,2,3$ forces, respectively,
$D_{12}=D_{23}=D_{31}=0$. These are all six possible edges, so $D=0$.

The map $P_v$ depends only on the local star and always has coordinate
zero equal to zero. On a cut $x=\delta s$ with $s_0=0$, it satisfies
$P_v(x)=s$. For an arbitrary $x$, the four messages are
\[
 (x_{01},\ x_{12}+x_{10},\ x_{23}+x_{20},\ x_{31}+x_{30}).
\]
Set $x_{02}=x_{23}=1$ and all other edge labels to zero. The transcript
is zero, as for the zero cut. But the triangle $(0,1,2)$ has label sum
one, so $x$ is not a cut. No referee receiving these messages can
recognize $\Cut_4$ exactly.
\end{proof}

This is an exact finite counterexample to the specified zero-root-row
template. It does not refute an asymptotic nonlinear lower bound or
show that every possible cut/noncut collision has been excluded.

\section{The cost of describing public graph lifts}

\begin{proposition}[Public-base reconstruction]\label{prop:lifts}
Fix a public simple base graph $H$ on $m$ labeled vertices and an integer
$r\ge2$. The input has public labels $(v,a)\in V(H)\times[r]$ and is
promised to be an $r$-sheet lift: every base edge is replaced by a
perfect matching between its two fibers, and there are no other edges.
There is an exact deterministic simultaneous protocol reconstructing
the entire input, and hence deciding connectivity, with at most
\[
 \Delta(H)\ceil{\log_2 r}
\]
bits per lifted vertex and polynomial decoding time.
\end{proposition}

\begin{proof}
Order the neighbors of each base vertex using the public labels.
Vertex $(v,a)$ has exactly one neighbor in each fiber $u$ with
$u\in N_H(v)$. In that order, it sends the sheet index of each such
neighbor, using $\ceil{\log_2 r}$ bits per index. Its message is
computed from its own neighborhood. The referee knows the order and
the base labels, so it reconstructs all edges from these indices.
Ordinary graph search decides connectivity. The encoding and decoding
are polynomial in the explicit input and output sizes.
\end{proof}

For $H=K_m$ and $n=mr$, the budget is
$(m-1)\ceil{\log_2 r}=O((n/r)\log r)$. Along a growing sequence of
such families, for any fixed $C\ge0$, the condition
\[
 \frac{r}{\log r}=\omega(\log^C n)
\]
makes this budget $o(n/\log^C n)$ and therefore rules out an
$\Omega(n/\log^C n)$ lower bound on that promise family.
The public base and public fiber labels are essential assumptions.
The proposition supplies no such protocol when these data are hidden,
and no upper bound for arbitrary graphs.

\clearpage
\section*{Provenance and open scope}

The baseline is the $k=1$ specialization of the supplied original draft,
\path{sources/draft/paper-body.pdf}, Section~6, pages 12--16.
Its explicit reconstruction is in
\path{results/agents/lower_bounds/2026-10-29-analysis.md}.
The four-symbol bound is proved here in full; the draft attributes its
general form to Karpovsky--Milman. No external citation is needed for
the arguments in this note.

The separate proof audit is recorded in
\path{results/agents/proof_auditor/2026-10-29-blind.md} and
\path{results/agents/proof_auditor/2026-10-29-audit.md}.
The blind report independently proved the cut-rank bound, the
four-vertex counterexample, the public-lift protocol, and a version of
Theorem~\ref{thm:barrier} with two extra output bits. It initially left
the central baseline unresolved; the subsequent audit checks the
released self-contained baseline proof.

Theorem~\ref{thm:barrier} uses the auditor's symplectic counting proof
and its subsequent refinement removing those two bits. The refinement
also received a separate reverse audit in
\path{results/agents/algebraic_sketches/2026-10-29-rank-constant-audit.md}.
The original algebraic analysis is
\path{results/agents/algebraic_sketches/2026-10-29-analysis.md}.
These records document independent checks within this research run;
they do not establish publication priority or author acceptance.

An unrestricted improvement beyond $\Omega(\sqrt n)$ remains open in
this work. General-$k$ claims, randomized protocols, stronger nonlinear
collision conjectures, and unselected prior-note assertions receive no
verdict here. This document is independently prepared; the original
author's editable source was not supplied, and the source PDF is not
modified.

\end{document}
